\documentclass[10pt,a4paper]{amsart}
\usepackage{amsmath,amssymb,amsthm,graphicx}
\usepackage{hyperref}
\usepackage{tikz}
%\usepackage{xcolor}
\usepackage{pict2e}
\let\arc\relax
%\usepackage{harpoon}
\usepackage{curves}
%\usepackage{chngcntr}
%\usepackage{apptools}
\usepackage{enumerate}

%\usepackage{verbatim}
%\usepackage[nospace,noadjust]{cite}
\usepackage{accents}
\usetikzlibrary{calc,decorations.markings}
\unitlength1cm

%%%%%%%%%%%%%%FONTS%%%%%%%%%%%%%%%%%%%%%%%%%%%%%%
\newcommand{\R}{\mathbb{R}}
\newcommand{\Z}{\mathbb{Z}}
\newcommand{\N}{\mathbb{N}}
\newcommand{\C}{\mathbb{C}}
\newcommand{\Id}{I}

%%%%%%%%%%%%%%MATHCAL%%%%%%%%%%%%%%%%%%%%%%%
\newcommand{\cK}{\mathcal{K}}
\newcommand{\cE}{\mathcal{E}}
\newcommand{\cR}{\mathcal{R}}
\newcommand{\cW}{\mathcal{W}}

%%%%%%%%%%%%%%%%%%ABBRS%%%%%%%%%%%%%%%%%%%%%%%%%%%%%

%%%%%%% OBJECTS SPECIFIC TO THIS PAPER%%%%%%%%%

\newcommand{\hw}{\widehat w}
\newcommand{\tw}{\widetilde w}

\newcommand{\bQ}{ Q}
\newcommand{\hQ}{{\widehat Q}}
\newcommand{\tQ}{{\widetilde Q}}

\newcommand{\sbQ}{Q^\star}
\newcommand{\shQ}{{\widehat Q^\star}}
\newcommand{\stQ}{{\widetilde Q^\star}}

\newcommand{\smu}{\mu^\star}
\newcommand{\shmu}{\widehat \mu^\star}
\newcommand{\stmu}{\widetilde \mu^\star}

\newcommand{\bmu}{{\mu}}
\newcommand{\hmu}{{\widehat \mu}}
\newcommand{\tmu}{{\widetilde \mu}}


\newcommand{\bP}{P}
\newcommand{\hP}{{\widehat P}}
\newcommand{\tP}{{\widetilde P}}

\newcommand{\sbP}{P^\star}
\newcommand{\shP}{{\widehat P^\star}}
\newcommand{\stP}{{\widetilde P^\star}}

\newcommand{\hR}{{\widehat R}}

\newcommand{\bv}{v}
\newcommand{\hv}{{\widehat v}}
\newcommand{\tv}{{\widetilde v}}

\newcommand{\sbv}{v^\star}
\newcommand{\shv}{{\widehat v^\star}}
\newcommand{\stv}{{\widetilde v^\star}}

\newcommand{\stu}{{\widetilde u^\star}}

\newcommand{\hF}{\widehat F}
\newcommand{\hW}{\widehat W}

\newcommand{\Fw}{\frac{F^2}{w}}
\newcommand{\hFw}{\frac{\hF^2}{\hw}}

\newcommand{\Cau}{\mathcal C_{\Sigma}}
\newcommand{\Cauv}{\mathcal C_{v}}
\newcommand{\Cauhv}{\mathcal C_{\hv}}
\newcommand{\Cautv}{\mathcal C_{\tv}}

\newcommand{\sCau}{\mathcal C_{\Sigma^\star}}
\newcommand{\sCauv}{\mathcal C_{\bv^\star}}
\newcommand{\sCauhv}{\mathcal C_{\hv^\star}}
\newcommand{\sCautv}{\mathcal C_{\stv}}

\newcommand{\Ltwo}{L^2(\Sigma)}
\newcommand{\sLtwo}{L^2(\Sigma^\star)}

\newcommand{\Sigl}{\Sigma^\ell}
\newcommand{\Sigr}{\Sigma^r}

%%%%%%%%%%%%%%%%OTHER SHORTCUTS%%%%%%%%%%%
\newcommand{\nn}{\nonumber}
\newcommand{\ul}{\underline}
\newcommand{\ol}{\overline}
\newcommand{\pa}{\partial}
\newcommand{\ti}{\tilde}
\newcommand{\wti}{\widetilde}
\newcommand{\wha}{\widehat}
\newcommand{\norm}[1]{\lVert#1 \rVert}
\newcommand{\abs}[1]{\lvert#1 \rvert}
%\newcommand{\Id}{\mathbb{I}}
\newcommand{\I}{\mathrm{i}}
\newcommand{\E}{\mathrm{e}}
\newcommand{\re}{\mathop{\mathrm{Re}}}
\newcommand{\im}{\mathop{\mathrm{Im}}}
\newcommand{\eps}{\varepsilon}
\newcommand{\dist}{\mathop{\mathrm{dist}}}

\newtheorem{Theorem}{Theorem}[section]
\newtheorem{Corollary}[Theorem]{Corollary}
\newtheorem{Lemma}[Theorem]{Lemma}
\newtheorem{Proposition}[Theorem]{Proposition}
 { \theoremstyle{definition}
\newtheorem{Definition}[Theorem]{Definition}
\newtheorem{Note}[Theorem]{Note}
\newtheorem{Example}[Theorem]{Example}
\newtheorem{Remark}[Theorem]{Remark} }

\numberwithin{equation}{section}
\usepackage{cleveref}
\crefname{originalequation}{original equation}{original equations}
\Crefname{originalequation}{Original equation}{Original equations}
\makeatletter
\newlabel{OuterParametrix}{{internal}{source}{}{sourcepointer-1}{}}
\newlabel{OuterParametrix@cref}{{[originalequation][0][]internal}{[1][0][]source}{}{}{}}
\makeatother
\crefname{originalequation}{original equation}{original equations}
\Crefname{originalequation}{Original equation}{Original equations}
\makeatletter
\newlabel{ParaConv}{{internal}{source}{}{sourcepointer-2}{}}
\newlabel{ParaConv@cref}{{[originalequation][0][]internal}{[1][0][]source}{}{}{}}
\makeatother
\crefname{originallemma}{original lemma}{original lemmas}
\Crefname{originallemma}{Original lemma}{Original lemmas}
\makeatletter
\newlabel{PropP}{{internal}{source}{}{sourcepointer-3}{}}
\newlabel{PropP@cref}{{[originallemma][0][]internal}{[1][0][]source}{}{}{}}
\makeatother
\title{Growth of the Legendre parametrix quotient near the endpoint plus one}
\author{Deift, Percy, Piorkowski, Mateusz}
\date{doi:10.3842/SIGMA.2024.004; CC BY 4.0}
\begin{document}
\maketitle
\setcounter{section}{5}
\setcounter{Theorem}{2}
\noindent\textit{Editorial note.} All delivered mathematical text is the author's own, copied from the pinned SIGMA article source; no statement or proof step was rewritten, reordered or removed. The counted statement is the corollary environment at source lines 1460-1466 and the counted proof is the proof environment at source lines 1467-1587. The environment itself carries no label; the anchor is the label of the first displayed alignment inside the statement, exactly as the intake census pairs it.\par\medskip
\noindent\textit{Editorial note.} The counted proof contains the author's own drawing of the contour, delivered verbatim as part of the proof text.\par\medskip
\noindent\textit{Editorial note.} The proof cites published references of the author's bibliography by key; the cited entries are copied line for line from the author's own bibliography into a retained bibliography fragment, and hyperref is added because the publisher class, not the author preamble, supplied the link command.\par\medskip
\noindent\textit{Editorial note.} The retained preamble of this unit was generated mechanically for this delivery: it copies the author's own package, numbering, macro and theorem declaration block from the pinned source (source lines 9-133), with the publisher class and page setup lines removed.\par\medskip
\noindent\textit{Allowed standard prerequisites (editorial).}\begin{itemize}
\item The construction of the local parametrix at the endpoint minus one. same article, the lemma at source lines 1363-1390. Provides the pattern that is copied at the endpoint plus one, as the author states at the beginning of the proof.
\item The outer parametrix and its matching with the local ones. same article, the display at source line 1351 and the matching condition at source line 1377. Supplies the normalisation with which the constructed local parametrix agrees up to an error of the order of the reciprocal of the parameter.
\end{itemize}\medskip
\ifdefined\hypertarget\hypertarget{sourcepointer-1}{}\fi
\noindent\textit{Source reference.} Original equation internal: sources/sigma-2024-004-recurrence-coefficients-for-orthogonal-polyn/pinned-original-source.tex, source lines 1351--1355. The counted statement or proof refers to that the author's displayed outer parametrix; it is the author's own earlier display, definition or statement in the same article and the whole argument of the counted span is retained verbatim.\par\medskip
\ifdefined\hypertarget\hypertarget{sourcepointer-2}{}\fi
\noindent\textit{Source reference.} Original equation internal: sources/sigma-2024-004-recurrence-coefficients-for-orthogonal-polyn/pinned-original-source.tex, source lines 1377--1379. The counted statement or proof refers to that the author's displayed matching relation between the quotients and the identity; it is the author's own earlier display, definition or statement in the same article and the whole argument of the counted span is retained verbatim.\par\medskip
\ifdefined\hypertarget\hypertarget{sourcepointer-3}{}\fi
\noindent\textit{Source reference.} Original lemma internal: sources/sigma-2024-004-recurrence-coefficients-for-orthogonal-polyn/pinned-original-source.tex, source lines 1363--1390. The counted statement or proof refers to that the author's lemma collecting the properties of the three local parametrices; it is the author's own earlier display, definition or statement in the same article and the whole argument of the counted span is retained verbatim.\par\medskip
\begin{Corollary}
For $z$ in a neighbourhood $U_{+1}$ of $+1$, the matrix-valued function $\tQ$ satisfies the asymptotics
\begin{align}\label{LegendreQBehaviour}
 \tQ(z) = O\big(|z-1|^{-1/4}\big)
\end{align}
uniformly for $n\to\infty$. Moreover, $\tQ$, and its boundary values $\tQ_\pm$ on $\Sigma$, are bounded in $\overline{\C} := \C \cup \lbrace \infty \rbrace$ away from small neighbourhoods of $\pm 1$, uniformly for $n\to\infty$.
\end{Corollary}

\begin{proof}
Analogously to the local parametrix for the Legendre problem at $z\! =\! -1$, ${\tP_{-1}(z)\! := \!\tP(z)}$, one can construct a local parametrix $\tP_{+1}(z)$ for the Legendre problem near the point $z = +1$ with local jumps inside an open neighbourhood $U_{+1}$ of $+1$ as depicted in Figure~\ref{contour1}.

\begin{figure}[th]
\centering\vspace*{3mm}
\begin{picture}(7,5.2)
\put(0.42,2.5){\line(1,0){3.08}}
%\put(5.5,2.5){\vector(1,0){0.2}}
\put(1.5,2.5){\vector(1,0){0.2}}


\put(1.35,0.35){\line(1,1){2.13}}
\put(2.1,1.1){\vector(1,1){0.2}}

\put(1.35,4.65){\line(1,-1){2.13}}
\put(2.1,3.9){\vector(1,-1){0.2}}


\put(2.3,3.9){$\big(\begin{smallmatrix}
1 & 0
\\
\phi^{-2n} & 1
\end{smallmatrix}\big)$}
\put(1,2.8){$\big(\begin{smallmatrix}
0 & 1
\\
-1 & 0
\end{smallmatrix}\big)$}
\put(2.3,0.9){$\big(\begin{smallmatrix}
1 & 0
\\
\phi^{-2n} & 1
\end{smallmatrix}\big)$}

\put(3.5,2.6){$+1$}
\put(4.5,3.5){$U_{+1}$}

\curvedashes{0.05,0.05}
\curve(3.4,2.5, 6.4,2.5)

\put(3.4,2.5){\circle{5.97}}
\end{picture}
\vspace{12pt}
\caption{The jump conditions for $\tP_{+1}$.}\label{contour1}
\end{figure}

In fact, we have $\tP_{+1}(z) = \sigma_3 \tP_{-1}(-z)\sigma_3$, see \cite[equation~B.10]{CD}. Hence, it follows from Lemma~\ref{PropP}\,(iv) that
\begin{align}\label{Legendre+1Esti0}
 \tP_{+1}(z) = O\big(|z-1|^{-1/4}\big)
\end{align}
uniformly for $n\to\infty$. After deforming the local contour such that it matches locally with $\Sigma$ as depicted in Figure~\ref{contour2}, we can analytically continue $\tP_{+1}$ as necessary to obtain a \emph{deformed} local parametrix $\tP_{+1}^{\text{def}}$, which would satisfy locally the same jump conditions as $\tQ$.

\begin{figure}[th]
\centering\vspace{2mm}
\begin{picture}(7,5.2)
\put(0.42,2.5){\line(1,0){4.58}}
\put(5, 2.5){\line(-1,0.59){3.65}}
\put(5, 2.5){\line(-1,-0.59){3.65}}
%\put(5.5,2.5){\vector(1,0){0.2}}
\put(2.8,3.8){\vector(1,-0.6){0.2}}
\put(2.8,1.2){\vector(1,0.6){0.2}}
\put(1.5,2.5){\vector(1,0){0.2}}
\put(4.4,2.5){\vector(-1,0){0.2}}

\put(2.3,4.3){$\big(\begin{smallmatrix}
1 & 0
\\
\phi^{-2n} & 1
\end{smallmatrix}\big)$}
\put(1,2.8){$\big(\begin{smallmatrix}
0 & 1
\\
-1 & 0
\end{smallmatrix}\big)$}
\put(2.3,0.5){$\big(\begin{smallmatrix}
1 & 0
\\
\phi^{-2n} & 1
\end{smallmatrix}\big)$}

\put(4.6,1.4){$\big(\begin{smallmatrix}
1 & 0
\\
2\phi^{-2n} & 1
\end{smallmatrix}\big)$}

\put(3.5,2.6){$+1$}
\put(3.4, 2.415){$\boldsymbol{\cdot}$}
\put(5,2.6){$1+\delta$}
\put(4.5,3.5){$U_{+1}$}

%\curve(1.35,0.35, 5,2.5)
%\curve(1.35,4.65, 5,2.5)

\curvedashes{0.05,0.05}
\curve(5,2.5, 6.4,2.5)
\curve(4.6,1.7, 4.1, 2, 4, 2.5)
\put(4,2.3){\vector(0,1){0.2}}

\put(3.4,2.5){\circle{5.97}}
\end{picture}
\vspace{15pt}
\caption{The jump conditions for $\tP_{+1}^\text{def}$.}\label{contour2}
\end{figure}

Because the jump matrices $\tv$ and their analytic continuations are uniformly bounded near $z = +1$, $\tP_{+1}^{\text{def}}$ would continue to satisfy the estimate \eqref{Legendre+1Esti0}
\begin{align}\label{Legendre+1Esti1}
 \tP_{+1}^{\text{def}}(z) = O\big(|z-1|^{-1/4}\big)
\end{align}
uniformly for $n\to\infty$. As the contour deformations are local, we also know that the matching condition \eqref{ParaConv} remains unchanged, at least on $\partial U_{+1}$:
\begin{align}\label{Legendre+1Esti2}
\tQ(s)\big[\tP_{+1}^{\text{def}}(s)\big]^{-1} = \Id + O\big(n^{-1}\big), \qquad s\in \partial U_{+1}.
\end{align}
However, as $\tQ\big[\tP_{+1}^{\text{def}}\big]^{-1}$ has no jumps inside $U_{+1}$, we can extend \eqref{Legendre+1Esti2} to all of $U_{+1}$ by the maximum modulus principle for holomorphic functions. Together with \eqref{Legendre+1Esti1}, the estimate~\eqref{LegendreQBehaviour} follows.

Regarding the second statement, it has been shown in \cite[Section~7]{KMVV} that for a wide class of weight functions including $\tw$, the outer parametrix solution $N$ is an approximation to the exact RH solution, in our case $\tQ$, uniformly in $z$ as long we stay away from the points $\pm 1$:
\begin{align*}
 \big|\tQ(z)-N(z)\big| = O\big(n^{-1}\big), \qquad z \mbox{ staying away from } \pm 1.
\end{align*}
As can be seen from \eqref{OuterParametrix}, $N$ is $n$-independent and bounded away from the points $\pm 1$. This finishes the proof.
\end{proof}

\begin{thebibliography}{99}
\bibitem{CD}
Conway T.O., Deift P., Asymptotics of polynomials orthogonal with respect to a
 logarithmic weight, \href{https://doi.org/10.3842/SIGMA.2018.056}{\textit{SIGMA}} \textbf{14} (2018), 056, 66~pages,
 \href{https://arxiv.org/abs/1711.01590}{arXiv:1711.01590}.


\bibitem{KMVV}
Kuijlaars A.B.J., McLaughlin K.T.-R., Van~Assche W., Vanlessen M., The
 {R}iemann--{H}ilbert approach to strong asymptotics for orthogonal
 polynomials on {$[-1,1]$}, \href{https://doi.org/10.1016/j.aim.2003.08.015}{\textit{Adv. Math.}} \textbf{188} (2004), 337--398,
 \href{https://arxiv.org/abs/math.CA/0111252}{arXiv:math.CA/0111252}.
\end{thebibliography}
\end{document}
